% =====================================================================
% METHODS
%
% FLAT HEADINGS, which is what Nature asks for: "subdivided by short bold
% headings referring to methods used", each at most 40 characters. There
% is no nesting and no numbering in the submitted document. The group
% names below (THE CORPUS, INPUT REPRESENTATION, ...) are LaTeX comments
% only -- they are how we navigate the file, not structure in the paper.
%
% WHAT THE GREY NOTES ARE FOR. Each one states what its passage DOES --
% the question it answers, what it establishes, and what the next passage
% needs from it. Read end to end they are the spine of the section, so
% the flow can be judged before any of it is written. Provenance and
% word counts come last in each note, not first.
%
% ARROWS between groups name the question being handed forward.
%
% WORD BUDGET: ~3,500 target. Allocation: corpus 500 . input 500 .
% encoder 450 . latent 150 . decoder 700 . objective 450 .
% evaluation 550 . reporting 200.
%
% THE SINGLE-BIOPSY CONFOUND IS A THREAD, NOT A SUBSECTION. Stated in the
% Introduction, then called back wherever it forces a choice:
%   CALLBACK 1  donor-level splitting      (Study selection)
%   CALLBACK 2  the adversarial head       (Decoder)
%   CALLBACK 3  cross-donor evaluation     (Evaluation)
% =====================================================================

\section{Methods}

The methods section elaborates on the full \elaj pipeline. Starting from 
the data and its representation, we deconstruct all major architectural 
choices made in the model, and finally close with optimisation and 
efficiency measures. The results, ablation studies, and discussions are 
separated as this section only details the hypotheses for the choices made 
and not the outcome of those choices.

% =====================================================================
%  THE CORPUS  -- budget 500.
% =====================================================================

\subsection*{Study selection}

The corpus used in this project is a complete record of publicly 
available single-cell data for endometriosis. It combines eleven
single-cell and single-nucleus studies, 1{,}577{,}728 cells, 254 donors
and six anatomical sites.

The search followed the PRISMA 2020 reporting guideline. Five public
archives were queried automatically: OmicsDI, NCBI GEO, BioStudies and
ArrayExpress, CELLxGENE Discover, and the Human Cell Atlas Data Portal.
The term \texttt{endometriosis} was matched against title, description
and keyword fields over the fifteen years preceding the search. A sixth
archive, CNGBdb, was searched by hand.

A record had to satisfy three conditions to enter the corpus: the data
had to be of or contain single-cell RNA sequencing information; the
sample had to be biopsied human tissue; and the data had to be
relevant to endometriosis. The last constraint was left understated to
accommodate relevant studies that are not necessarily of endometriosis
patients. These studies were either of adenomyosis or of a normal 
physiological state. Adenomyosis was added because it is 
pathophysiologically related to endometriosis. Physiologically stable 
tissue was included to augment the control samples from other studies.

The corpus was de-duplicated by donor by referencing the paper published
for each dataset. Each donor was given a unique identifier. Tissue site
was mapped to UBERON terms and disease state to MONDO terms, so that a 
cell's origin is recorded in a standard form other datasets in the 
corpus can follow. No correction was applied to the expression values 
themselves.

Every study identifies the sample each cell came from, so every cell can
be traced to a donor. The corpus was split into training, validation and
test partitions targeting a 70/15/15 ratio, with every cell from a donor
placed entirely in one partition. Cells from one donor share a genotype,
a hormonal state, a handling protocol and in most cases a single biopsy,
so a cell-level split would place near-duplicates on both sides and
measure the model on cells it had effectively already seen. One study
pair deposits the same fifteen people under two sets of identifiers,
which the split could not detect; 17{,}873 held-out cells, 4.3\% of the
held-out set, therefore share a donor with training.


\todo{PRISMA figure. Its own display item, or panel (a) of a composite
Figure 1. Nature counts items not panels, so a composite costs nothing in
the item budget and only costs text length.}


\subsection*{Gene vocabulary}

The datasets in this corpus cover a heterogeneous set of genes. Gene
names are also not stable.

In an effort to standardise, every gene was resolved against the HGNC 
approved symbol set, with anything unresolved dropped and recorded against 
its study of origin. An alias pointing to two different approved symbols 
is ambiguous, and resolving it would mean guessing; all such aliases were 
dropped, removing 1{,}452 mappings.

With the names settled, the set of genes the model reads still has to be
chosen. The selection used here is derived from Reactome or Gene Ontology 
(GO) pathway groups.  A fixed vocabulary of 15{,}496 genes was taken as those 
annotated to at least one Reactome or GO pathway, on the reasoning that a
gene worth modelling is one known to participate in a biological process. 
This has added benefits apparent in the encoder later. Restricted to this 
vocabulary, the proportion of genes a pair of studies holds in common 
rises from 64–96\% to 87–98\%.

\medskip\noindent\textcolor{elajmuted}{$\rightarrow$ \textit{The corpus is
one object with a chosen vocabulary. How does a cell's counts over that
vocabulary become something a transformer can read?}}

% =====================================================================
%  INPUT REPRESENTATION  -- budget 500. NOT YET SEEDED.
% =====================================================================


\subsection*{Gene tokens}

A cell is passed to the model as a single datapoint: a sequence of
tokens, one per gene in the fixed vocabulary.

Each token is the sum of two parts. The first is a learned identity
embedding ($e_i$). This is a vector unique to a gene fixed for all cells. 
The objective of this embedding is to encode which gene we are referring 
to. It is a $d_{\text{token}}$ dimensional vector.

The second part encodes the magnitude of a gene's expression in a cell 
($c_i$). Gene expression in the raw input data is a number representing 
the occurrence of a gene in a given cell. This one-dimensional value is 
first library-normalised to a target sum of $10^4$ and log-transformed. 
The transformed value still needs to be projected to $d_{\text{token}}$
dimensions from one. To project it without collapsing all genes onto a 
single direction, we multiply the transformed value by a learnable 
square matrix of size $d_{\text{token}}$ ($A$) which is shared across 
all genes. The resultant scaled matrix ($c_i A$) is then multiplied by 
the gene identity vector to get a $d_{\text{token}}$ dimensional vector 
for gene expression. A gene token is therefore computed as:

\begin{equation}
\text{token}_i \;=\; e_i + c_i A e_i .
\label{eq:token}
\end{equation}

By using the gene identity vector in projecting gene expression values
we are able to avoid a trivial low-dimensional projection. The directions
of each gene expression value naturally differ because the gene identity
embedding is unique per gene. A trivial projection would collapse a
cell's entire expression profile onto a single line. Conversely, a 
per-gene learned expression embedding would also allow unique gene 
directions but introduce roughly 500{,}000 parameters to the model. 
Through this method we get away with a single matrix of roughly 1{,}000 
parameters and the identity and expression stay separable by construction.


\subsection*{Unmeasured genes and the sentinel}

Every entry in the corpus is marked as one of three things: a count, a 
measured zero, or not measured in this study. The third case is stored 
explicitly because zero is a value of expression and non-measured genes
cannot be assumed to have a value. None of the models compared in this paper 
carries an input state for a gene that was not assayed. While it can be 
argued that most entires in the raw counts are zero and thus it would 
make little difference to the model, we found it important to specify 
unmeasured genes to the model as a Did Not Sample (DNS) sentinel.

DNS sentinel is a distinct and reserved learned vector in place of a 
value embedding. Crucially, it does not impact the identity embedding of a 
gene. Which means it doesn't interfere with what a gene is, only how much it 
is measured. Unmeasured genes are also excluded from the reconstruction loss, 
since there is no observation to score a prediction against. This allows the 
model to develop a holistic and true understanding of a cell. More details 
about how much it affects the outcome are in the ablation studies section 
under results.


\subsection*{Pathway gene grouping}

Self-attention compares every token with every other token, so its cost
grows with the square of the sequence length. One token per gene over
15{,}496 genes is roughly 240 million pairs for a single cell. That is
where the parameters and the compute of most single-cell transformers go,
and it is why those models need to be trained on tens of millions of cells. 
This corpus has 1.58 million cells.

Previous models have either reduced the number of genes studied, compressed 
their input representation in a lossy manner, or just used a very large 
corpus. None of these are transferable to this model. Reducing the number 
of genes to the top two or three thousand most variable genes assumes that the 
genes that vary the most are the most relevant to the disease. We do not have 
any reason to believe this assumption holds for endometriosis. A lossy 
compression for input representation comes with similar assumptions.

Instead, we used the partitioning available through the Reactome and GO 
hierarchies and cut them into 44 groups. The hypothesis being that genes 
that work together biologically would have more interesting interactions 
and interactions between groups of genes would be more relevant than 
between individual genes belonging to different groups. Following this idea, 
attention is first run within each group and attention across groups is 
run afterwards. Consequently, attention work falls from 240 million pairs 
to 5.72 million, a reduction of 97.6\%, and no data-losing assumptions 
have been made.

Apart from parameter efficiency, gene grouping also increases the effective
subspace in which gene identity lives. Each group is processed independently, and
no cross-group gene interactions are attended to. The gene space is 
therefore block diagonal, one block per group, and the dimension actually 
available to represent a gene is the token width times the number of groups,
$32 \times 44 = 1{,}408$, rather than the token width alone. For
comparison, Gene2vec, which is the conventional off-the-shelf gene
embedding, uses 512 dimensions. Effectively, a 32-dimensional token here
is a slice of a gene space nearly three times wider than the standard one. The 
derivation is given in Supplementary Note 1.

The reason it is an effective space and the reason it is able to encode 
more information is the same: the block diagonal nature of the gene space 
matrix. Since a block diagonal matrix is a special case of a regular matrix, 
what each block contains is critical to expressing all genes. That is why 
the pathway grouping is relevant. Without that, we would be restricting the 
gene space to block diagonal and getting nothing in return. This is tested 
through an ablation study.


\medskip\noindent\textcolor{elajmuted}{$\rightarrow$ \textit{A cell is now
44 groups of tokens. What reads a group, and what does it hand on?}}

% =====================================================================
%  ENCODER  -- budget 450. Needs MORE than the thesis has: the funnel,
%  and how information is squeezed.
% =====================================================================

\subsection*{Group encoders}

Each group is read by its own transformer: two layers of multi-head
self-attention, four heads, operating on that group's input tokens at width
32. Within a group, attention models what any two genes are doing 
together in a cell. The summary of interactions inside each group is 
collected in a CLS token. What leaves the stage is essentially a summary of 
each group. 

Notably, this is the first and most stark compression in the model. A cell 
enters as 15{,}496 gene tokens and 44 summary tokens are created. A 350-fold 
reduction. Additionally, a learned CLS token reads the group through a single
query, and what it carries away is a weighted combination of what
that one query attended to. Which is to say, it is a lossy form of compression. A 
weighted average discards the shape of a distribution. To recover more 
information, \elaj{} concatenates eight order statistics computed over the 
group's gene tokens to the learned token. The mean, the standard deviation, 
the minimum, the three quartiles, the maximum, and the mean of the top four 
are the statistics appended. Each statistic becomes its own summary token 
rather than being folded into one wide vector, so the next stage sees them 
separately.

A group therefore hands on nine summary tokens instead of one, and the
hypothesis being made is explicit: a single learned query is a lossy
readout of a group of genes, and that order statistics recover part of
what the averaging throws away. Ablation C-4 removes the eight statistics
and keeps the learned token alone, taking what each group contributes from
288 numbers back to 32.

Now, a cell enters as 15{,}496 gene tokens and leaves this stage as 396 
summary tokens. A reduction of roughly 39-fold achieved without attending 
across the vocabulary and without dropping a gene.


\subsection*{Cell aggregation}

Up to this point the model has only looked at genes within different 
pathways independently. The second stage asks how the groups relate to one 
another.

The 396 group summaries are assembled into one sequence together with a
learnable cell summary token appended. The second stage of the encoder 
is a transformer with two layers, four heads, and width 64. The learned 
cell summary token is the first representation of the complete cell 
inside the model.

The purpose of this stage is to allow the model to understand interaction 
between the individually studied gene groups. The hypothesis being that 
at this stage the model learns biologically relevant cellular information 
from an analysis of known pathway groups.

The lossy compression objection raised one level down applies here too. 
Consequently, the same remedy is applied again: order statistics are 
computed over the group-token positions and merged with the cell token. 
The merge projects back to width 64 rather than widening the output.

A cell therefore leaves the encoder as a single vector of width 64,
narrowed from 15{,}496 gene tokens through 396 group summaries.


\medskip\noindent\textcolor{elajmuted}{$\rightarrow$ \textit{A cell is one
vector of width 64. What does the model keep when it is forced to describe
that cell in sixteen numbers?}}

% =====================================================================
%  THE LATENT  -- budget 150.
% =====================================================================

\subsection*{Latent space}

The latent is where the model's most compressed understanding of a cell 
is calculated. Everything before it works to populate it, and everything 
after it exists to check that what was put there is worth having. \elaj 
being a CVAE, its latent space is the main component of interpretation.

As with most CVAEs, \elaj compresses information into a Gaussian subspace. 
In practice, it is sixteen Gaussian distributions: sixteen means and 
sixteen variances. Compared to baselines, a sixteen dimensional subspace 
is tiny. We arrived at this experimentally rather than theoretically. More 
details are given in the results section.


\subsection*{Decoder conditioning}

Being a CVAE, \elaj passes metadata about a cell to the decoder. Only the 
deocder is conditioned in this model. Most importantly, the study a cell 
belongs to and the platform used for sequencing that cell, both of which 
help the model identify and isolate non-biological information, are 
passed as learned embeddings. Tissue is encoded into four levels: 
compartment, organ, tissue type, microsite. The combined metadata 
information is appended to the 64-dimensional projection of the 
16-dimensional latent space to give a combined 112-dimensional vector 
per cell to the decoder.

This whole exercise is done to remove batch effects and non-biologically 
relevant information from the latent. The exception is tissue details.
Tissue, referring to the site of biopsy, is provided to the decoder to 
remove the most obvious grouping of cells. Endometriosis is more than 
tissue identification. To uncover the biological signal beyond tissue, 
we remove the need to model tissue information.

Crucially, disease information is not given to the decoder. A decoder 
told which cells are diseased will not rely on the latent for the signal 
that separates a normal cell from a diseased one, which is 
counterproductive to understanding the through the latent space. Every 
disease result in this paper is therefore structure the model was never 
shown. However, one study is the sole source of both adenomyosis and 
myometrium in this corpus, so for those 49{,}606 cells the study 
embedding is also a disease indicator. Claims about adenomyosis are 
therefore selectively made.


\medskip\noindent\textcolor{elajmuted}{$\rightarrow$ \textit{The decoder
knows where the cell came from. Now it has to produce 15{,}496 numbers from
sixteen.}}

\subsection*{Reconstruction head}

The decoder of \elaj consists of 4 heads. The first, and the model’s 
primary objective, is reconstruction.

The reconstruction head turns the 112-dimensional context into one
predicted expression level for every gene the cell's study measures.
Between 13{,}144 and 15{,}056 genes are predicted, depending on the
study's DNS count. An output layer holding one row of weights per gene
would do this in 1{,}007{,}240 parameters, and would learn each gene's
response in isolation. \elaj scores genes against the cell instead, in
122{,}825 parameters.

\textbf{The design.} Each gene is given a 64-dimensional query vector,
derived from the 32-dimensional identity embedding the encoder uses as
input. The projections are unique per pathway group, making the 
weight and bias an accumulation of common information among genes in 
each group. This is what allows two genes in different groups to have 
similar identity embeddings without being confused for one another.

The cell context is read into two things: a 64-dimensional vector that 
scores against the gene query, and a scalar that scales the resulting 
profile. A gene's predicted log-expression is then given as:
\begin{equation*}
\log \mu_{ig} \;=\; \frac{(W_k c_i) \cdot q_g}{\sqrt{64}}
\;\big(\mathbf{w}^\top W_v c_i\big) \;+\; b ,
\end{equation*}
where $q_g$ is gene $g$'s query, $c_i$ is cell $i$'s context, and $W_k$,
$W_v$, $\mathbf{w}$, and $b$ are learned. $(\mathbf{w}^\top W_v c_i)$ is a
scalar value. Since there is no per-gene value vector of the kind a regular
attention layer would have, this is a bilinear scoring layer, closest in
form to scGPT's masked-value-prediction head.

\textbf{No per-gene baseline.} Every term above is multiplied by something
belonging to the cell, so no parameter gives a gene a constant expression
level. Most count decoders end in a per-gene linear layer whose bias
supplies exactly that, making the average expression profile free; here it
has to be delivered through the latent or the covariates.

\textbf{From scores to counts.} Capture depth varies by an order of 
magnitude across this corpus. So instead of asking the model to recreate 
the exact counts, we pose a simpler problem: predict how much each gene 
contributes to the total count of the cell. The model gets library size, 
total count of genes in a cell, as input for this task.

\textbf{The likelihood and the target.} The model predicts the mean 
count of each gene, and the observed counts are scored against a negative 
binomial with one dispersion learned per gene. The loss is the mean 
negative log-likelihood over all cells and all measured genes, so every 
reconstruction figure in this paper is in nats per gene. Genes hidden 
from the encoder during training are still demanded of the decoder, 
which makes this masked autoencoding.


\medskip\noindent\textcolor{elajmuted}{$\rightarrow$ \textit{A reconstruction
exists. What stops it being produced WITHOUT the latent?}}

\subsection*{Adversarial study removal head}

The second head of the decoder is an adversarial network designed to 
actively remove batch information from the latent. The network is trained to 
predict the study that a cell came from given the latent. It is given the 
latent means rather than sampling from the gaussian space. This 
removes variational noise that the adversarial head does not need to model. 
To make it an adversarial network, during backpropagation the gradient is 
reversed before reaching the latent, making discriminator performance a 
pressure against batch information.

Two weights govern this. One is the share of total loss attributed to the 
discriminator which governs how much the model optimises for discrimination. 
The other is the scaling factor of the reversed gradient. This determines 
how much the discriminator affects the encoder. It is the second that 
traces the frontier reported in this paper. 

As with any adversarial network, the discriminator's capacity is a sensitive 
hyperparameter. A strong discriminator prevents the model from learning 
anything while a weaker one is easily fooled by the encoder. The failure 
mode of a weaker discriminator is described in results. \elaj uses a 
discriminator with three linear layers of width 256, with a leaky rectifier, no
dropout, and layer normalisation after each layer.


\medskip\noindent\textcolor{elajmuted}{$\rightarrow$ \textit{One term pushes
structure OUT of the latent. The next two decide what stays in.}}

\subsection*{Elastic cell similarity head}

The third head of the decoder is a selective pressure head adapted from 
scGPT. Notably, this head acts pairwise on the cell embeddings within a 
batch and not against one cell. Every pair of cells is checked against a 
cosine similarity threshold of 0.8. Pairs falling below the threshold are 
pushed away from each other while cells that clear the threshold are pulled 
closer together and pairs at or below zero cosine are left untouched. 

In practice this acts as a sharpening term and its purpose is to make a 
collapsed latent expensive. A single undifferentiated cloud is as 
uninformative as clear study-sectioned clouds. This head stretches 
dissimilar cells in a batch while simultaneously pushing similar cells 
together. Crucially, every pair is drawn from one study because a batch is 
composed of one study by construction. A cross-study pair of the same cell 
type with a low cosine because of batch effects would be driven further 
apart instead of closer. That puts unnecessary pressure on the adversarial 
head. Confined to one study the term can sharpen within-study structure, 
which helps gather biologically similar cells together.


\medskip\noindent\textcolor{elajmuted}{$\rightarrow$ \textit{Both terms so
far shape the latent without being told anything. One term is told.}}

\subsection*{Menstrual-phase head}

The fourth head is the only head other than reconstruction that introduces
new information in the model rather than moulding existing information. 
Cycle phase is among the largest sources of biological variation in
endometrium. The menstrual phase head is a linear classifier 
that reads the Gaussian latent means and predicts the cell's 
menstrual-cycle phase.

The reconstruction head has no pressure to isolate menstrual cycle. Each 
donor contributes a single biopsy in this corpus which is a data confound 
that needs to be taken into consideration. That relates directly to the
adversarial head as it acts against menstrual cycle information simply 
because it has the same geometric form as a study effect. The menstrual 
phase head gives that trade-off a second side, so moving the adversarial 
strength traces a frontier rather than erasing both at once.

However, a head with hidden layers would learn how to separate the latent
nonlinearly which doesn't necessarily make the latent more informative. 
Instead \elaj uses a single linear map that puts the separation pressure 
directly on the latent, creating a more interpretable subspace. 

Notably, \elaj is only semi-supervised on menstrual phase because 
512{,}101 of 1{,}161{,}472 training cells have menstrual phase metadata.
Unlabelled cells are not penalised. The effects of this head are quantified
through ablation studies described in results.

Menstrual phase is the only label the model is given. Cell type, disease 
status, stage and grade have no head in any configuration reported here, 
so every result on those is unsupervised.


\medskip\noindent\textcolor{elajmuted}{$\rightarrow$ \textit{Five terms
exist. What does the sum look like?}}

\subsection*{Composite objective}

The four heads are trained together through one loss value. 
Reconstruction acts over all genes with weight 1 while every other 
term is weighted against it as:
\begin{equation*}
\mathcal{L} \;=\; \mathcal{L}_{\text{nb}}
\;+\; \frac{1}{G}\Big[\,
\beta(t)\, w_{\text{kl}}\, \mathcal{L}_{\text{kl}}
\;+\; w_{\text{ecs}}\, \mathcal{L}_{\text{ecs}}
\;+\; w_{\text{dab}}\, \mathcal{L}_{\text{dab}}
\;+\; w_{\text{cls}}\, \mathcal{L}_{\text{cls}}
\,\Big] ,
\end{equation*}
with $w_{\text{kl}}$ 1, $w_{\text{ecs}}$ 50, $w_{\text{dab}}$ 0.5 and
$w_{\text{cls}}$ 40. Reconstruction is a mean over genes while the four
bracketed terms are per-cell quantities, so each is divided by the number of 
genes that batch's study measured ($G$).

The price on the latent is the usual one: the Kullback--Leibler divergence
between each cell's posterior and a unit Gaussian, summed over the sixteen
dimensions and averaged over cells. $\beta(t)$ ramps it linearly from zero
over the first 5\% of training steps, reaching 0.5, because until
reconstruction is established the latent holds nothing worth paying for and
a full price collapses it first. $w_{\text{kl}}$ is the second of the
paper's two knobs: at a fixed adversarial strength it moves the model along
the frontier rather than displacing the frontier itself.
information cost and is described next.

$\lambda_{\text{dab}}$ is absent from this sum, which is deliberate. It
scales the gradient as it is reversed on the way to the encoder, so it
governs how much of the adversarial loss the latent feels while leaving the
discriminator's own training untouched. That is the knob which traces the
frontier, and it is set to 168 in the reported arm. Only the first two 
terms are the model's own objective. Checkpoint selection and early 
stopping read reconstruction alone.


\subsection*{Cross-donor agreement}

In this corpus, every donor was biopsied only once. This has repercussions in
how the model performance should be measured on this corpus. Because a donor
contributes a single sample, that donor has exactly one menstrual phase, one
disease status, one stage and one grade, and every cell drawn from them
inherits those values unchanged. Each of these labels is therefore a property
of the donor rather than of the cell. Cell type is the only exception: it
varies within a single biopsy.

The obvious way to ask whether a latent carries a label is to check whether
cells agree with their neighbours: take each cell's nearest neighbours in the
latent and measure how often they share its label. On this corpus that test
cannot be trusted, and the reason is the paragraph above. If a label belongs
to the donor, then asking whether a cell's neighbours share its label is
asking whether they come from donors carrying that label, and a cell's own
donor carries it by definition. Every neighbour drawn from the same donor is
an automatic match, scored whether or not the latent encodes anything about
the label at all.

Plain agreement therefore rests on a floor set by how many of a cell's
neighbours come from its own donor, and on this corpus that floor is high.
Among the hundred nearest neighbours of a labelled cell, between a quarter
and three-fifths belong to the same donor depending on the label: 37\% for
menstrual phase, 29\% for disease status, 61\% for stage. A latent carrying
no donor information at all would give the chance rate, the probability that
two cells drawn at random share a donor, which is between 0.6\% and 3.1\%
here. Same-donor neighbours are thus twenty to fifty times over-represented,
and because each is an automatic match that enrichment passes straight into
the score: of the agreement a plain measurement reports, between 44\% and
71\% is the free same-donor part.

Worse, it is a floor the model can raise. A model that learns to recognise
donors increases the share of its own donor among each cell's neighbours and
improves its own score, so ranking models by plain agreement partly ranks
them by how well they memorised the donor, the opposite of what this
paper is measuring.

\elaj\ is measured with the same-donor neighbours removed. For each label,
the hundred nearest neighbours of every labelled cell are found in the
sixteen-dimensional latent and collected into a donor-by-donor matrix: entry
$(d, d')$ counts how many neighbours of donor $d$'s cells belong to donor
$d'$. The diagonal is then discarded (that discarding
is what cross-donor means), and agreement is the share of the remaining neighbour relations
that join two donors carrying the same label. Chance is the probability that
two such relations drawn at random would share a label, $\sum_j p_j^2$ over
the label's shares across donors, and agreement is reported as excess over
it, $(a - c)/(1 - c)$, so labels with different numbers of levels are
comparable.

The null is obtained by reassigning labels to donors. Two hundred random
permutations of the donor-to-label map are scored through the same
expression, which asks what agreement this latent would produce if the same
labels had been dealt to different donors, the appropriate null for a
quantity constant within a donor. The reported $p$ is the share of
permutations reaching the observed agreement, counted as
$(1 + \#\{\text{null} \geq \text{observed}\}) / (1 + 200)$.

That construction puts a floor under $p$, and the floor matters for reading
this paper's negative results. With two hundred permutations the smallest
attainable value is $1/201$, and more restrictively the number of distinct
permutations is set by how many donors carry the label at all. Where a label
reaches only a handful of donors, no arrangement of them can produce a small
$p$, so a result that is not significant means the test could not detect an
effect of any size rather than that no effect exists. Disease stage and
endometriosis grade reach 58 and 60 donors respectively, against 269 for
disease status, and are reported as undetectable rather than absent.

Because the labels are constant within donors, the matrix is a sufficient
statistic: agreement and its null are arithmetic on a table of integers with
at most one row per donor, and the neighbour search that produced it is the
only step needing the whole corpus in memory. Any later regrouping of a label
is recomputed from the stored matrix without re-embedding a cell.

The construction is a permutation test of assortativity on a fixed graph.
The statistic is the share of edges joining same-label nodes, which is the
enumerative mixing coefficient of \citet{newman2003mixing} computed on a
donor graph rather than its usual node graph; significance by permuting group
identities against an unchanged proximity structure is the procedure of
\citet{anderson2001permanova}; and permuting at the level of the donor rather
than the cell is what \citet{squair2021confronting} prescribe, since donors
and not cells are the independent replicates in a multi-donor single-cell
study.

\note{THIS SUBSECTION IS NOT REVIEWED. AW, 2 November 2026. Everything
below was established by reading
\texttt{endo\_model/evaluation/sufficient\_stats.py} on 2 November and is
recorded here because the prose depends on it and because one error had
already propagated out of it. The same material is in
\texttt{axw1008/docs/paper/CROSS\_DONOR\_ESTIMATOR.md}.
\\
\\\textbf{WHAT THE TEST IS, IN ITS OWN VARIABLES.}
\\\quad $\bullet$ \texttt{matrix} is donor $\times$ donor. Entry $(d,d')$
counts how many of donor $d$'s cells' hundred nearest neighbours belong to
donor $d'$.
\\\quad $\bullet$ \texttt{np.fill\_diagonal(matrix, 0.0)} deletes the
same-donor relations. \textbf{The off-diagonal is the cross-donor part};
zeroing the diagonal is what makes it so. An earlier draft of this note had
that backwards.
\\\quad $\bullet$ \texttt{lab} holds \emph{one label per donor}, which is
only legitimate because the label is donor-constant.
\\\quad $\bullet$ \texttt{agree(assignment)} is
\texttt{(m * same).sum() / total}: of every cross-donor neighbour relation,
the share joining two donors with the same label.
\\\quad $\bullet$ \texttt{observed = agree(lab)}.
\\
\\\textbf{THE HYPOTHESES.} $H_0$: the latent's donor-level geometry is
unrelated to the label, so the observed agreement is what any assignment of
these labels to these donors would give. $H_1$: donors sharing a label sit
closer, so their cells fall in each other's neighbourhoods more often.
\\
\\\textbf{WHAT IS PERMUTED, AND WHAT IS NOT.} The embedding and the
neighbour graph are computed once and never change. Cells do not move and
donor membership does not change. Only \texttt{lab} is permuted: the same
labels are dealt to the same donors in a different order. The donor is the
exchangeable unit because the label is constant within one, so nothing below
donor level can carry information about it; permuting cell labels would break
that constancy and build a null so weak that everything would pass.
\\
\\\textbf{\texttt{chance} AND THE NULL ARE NOT THE SAME THING, AND ONLY ONE
IS CALIBRATED.} \texttt{chance} is $\sum_j p_j^2$ over the label's shares
\emph{across donors}, weighting every donor equally. But \texttt{observed}
weights each donor pair by the neighbour relations between them, so large and
tightly-clustered donors dominate it. The two are computed under different
weightings, which makes \texttt{excess} $= (a-c)/(1-c)$ a descriptive effect
size with an approximate correction rather than a calibrated one.
\textbf{The permutation $p$ is the inferential statistic}, because it
recomputes under the same weighting.
\\
\\\textcolor{elajcoral}{\textbf{LOOPHOLES, WORST FIRST. These belong in the
Discussion's limitations, not here.}
\\1. \textbf{\texttt{excess} is not comparable across labels, which was the
stated purpose of normalising it.} The null's own 95th percentile, carried in
the benchmark \texttt{note} column, tracks donor count: 0.06 at 34 donors
(cell family), 0.18 at 36 (phase), 0.12 at 18 (grade), \textbf{0.42 at 16}
(stage). Reassigning sixteen donors' labels at random reaches 0.42, so
reading 0.1313 and 0.7346 on one scale is wrong.
\\2. \textbf{$p$ is one-sided} --- \texttt{(null >= observed).sum()}. It can
never establish a significantly \emph{negative} value, so grade's
$-$0.0414 is ``not positive'', never ``negative''.
\\3. \textbf{For a label spanning several studies the permutation is too
permissive.} Disease status covers 269 donors across eleven studies and some
cohorts are entirely cases, so permuting across all of them destroys the
label--study association and the null never sees the study structure partly
producing the observed agreement. This is visible in our own numbers: pooled
disease status is significant at $p=0.005$, and inside each cohort, where the
permutation is necessarily within-study, nothing is
($p=0.3433$, $0.0697$, $1.000$). Stage and grade are unaffected because every
labelled donor for each sits in one cohort.
\\4. \textbf{The permutation is unstratified by donor size}, so one draw can
hand a label to donors carrying very different adjacency mass. A size-matched
permutation would give a tighter null.
\\5. $k=100$ is a free choice the statistic depends on, with no sensitivity
analysis.
\\6. $p$ floors at $1/201$, so 0.005 means ``nothing in two hundred beat
it'', not a measured small probability.}
\\
\\\textcolor{elajcoral}{\textbf{A CLAIM IN THE PROSE ABOVE IS WRONG AND
NEEDS CUTTING ON REVIEW.} ``the number of distinct permutations is set by how
many donors carry the label at all'' is not the binding constraint at 16 to
60 donors, where the number of distinct assignments is astronomical. What
binds is the \emph{width} of the null, which is loophole 1. The p-floor
argument should be made that way instead, and it is the stronger version.}}





\subsection*{Donor and study recovery}
\note{PURPOSE: define the two costs, so the frontier has both its axes
before it is plotted. These are the quantities a good model should fail
at.\\
HANDS ON: both axes exist. What does the latent's shape look like under
them?\\
SOURCE: Thesis 05 §Study mixing.}



\subsection*{Whole-corpus embedding}
\note{PURPOSE: licence the paper's discovery half. Once the model is
fixed, nothing is being selected, and the metadata read below was never
an input --- so there is nothing a held-out split protects against for
these labels, and all 1{,}577{,}728 cells can be embedded and read
together. Phase is the exception, governed by the head. Also the region
procedure: every cell takes the cell family of its 25 nearest LABELLED
neighbours searched in the sixteen dimensions, not in the UMAP, because a
bounding box assumes each family occupies one convex patch and immune
cells do not.\\
SOURCE: new.}

\medskip\noindent\textcolor{elajmuted}{$\rightarrow$ \textit{Everything
measurable has been measured. What has to be declared?}}

\subsection*{Optimisation}
\note{PURPOSE: how a run is executed, kept short. AdamW, $10^{-3}$, batch
256, bf16, cosine over a 20-epoch horizon chosen for its SHAPE rather
than its length --- the horizon sets how long the learning rate stays
high, and early stopping is what ends a run. Per-study batch construction
and domain-specific normalisation are two sentences here rather than
headings of their own, per your note that most training-procedure
material is statements belonging inside other passages.\\
HANDS ON: runs finish. Which checkpoint is the model?\\
SOURCE: Thesis §Configuration (72). Hyperparameter detail $\rightarrow$
Supplementary.}

% =====================================================================
%  REPORTING  -- budget 200. NOT YET SEEDED.
% =====================================================================



\subsection*{Use of large language models}
\note{PURPOSE: \textbf{REQUIRED BY NMI} --- ``Use of an LLM should be
properly documented in the Methods section.'' What was used, for what,
and what it did not do.\\
SOURCE: Thesis Appendix A, rewritten to that requirement.}

\subsection*{Data availability}

This study analysed previously published data and generated no new primary
data. All datasets are publicly available: Gene Expression Omnibus
accessions GSE266265, GSE247695, GSE214411, GSE218044, GSE213216, GSE203191
and GSE179640; ArrayExpress accession E-MTAB-14039; and CellxGene Discover
collections \texttt{9ddea8d9} and \texttt{968885c5}.
CNGBdb accession CNP0006183 (li et al.).


\subsection*{Code availability}

All code for data preparation, model training and evaluation is available at
\textcolor{elajcoral}{[repository URL]} under an
\textcolor{elajcoral}{[MIT]} licence.

